\chaptertitlepage{Conclusions}{Welfare regimes, late-life partitions, and the silver-economy design gap}
\label{ch:ch9}

\section{Overview}
\label{sec:ch9-overview}

This thesis has developed a multidimensional segmentation of the over-65 population in Italy and Sweden using SHARE Wave~9 data, with a methodological commitment to including subjective dimensions of late-life experience as constituent inputs to the K-means partition rather than as external validation outcomes. The empirical core, comprising Chapters~\ref{ch:ch4}--\ref{ch:ch7}, has produced two parallel typologies of five Italian and six Swedish archetypes, a formal Hungarian-matched cross-country comparison that identifies four matched pairs and three regime-specific unmatched profiles, three structural contrasts that distinguish the two typologies on the shape of fragility, on the operational meaning of digital adoption, and on the late-life trajectory of the high-resource pole, and a robustness battery that documents the structural persistence of the partition under stricter modelling alternatives and under the pre-pandemic Wave~8 reference. This chapter draws together the substantive contributions, develops the theoretical and operational implications, documents the boundaries of the analysis, and identifies the directions of future research that the empirical findings open.

\section{Substantive findings}
\label{sec:ch9-findings}

The empirical analysis has produced four findings that the thesis takes as its substantive contribution.

The first is the typological identification of the Isolation Paradox in the Italian over-65 population. The Moderate Isolated archetype (\S\ref{sec:ch4-moderate-isolated}, 26.9\% of the analytical sample) combines intact functional health and economic resources with thin social and digital engagement, and under-utilises the discretionary, navigation-intensive segments of the healthcare system: relative to the structurally comparable Traditional Social cluster, the regression-adjusted analysis of Table~\ref{tab:ch4-isolation-regression} returns 0.67 fewer specialist contacts per year and 27.9 percentage points lower preventive dental-care attendance (both $p < 0.001$). The Hungarian-matched analysis confirms the configuration is structurally Italian-only (\S\ref{sec:ch6-isolation-paradox-comparative}): it is the empirical signature of a familistic regime that offers no substitute infrastructure for individuals not embedded in dense informal networks, with under-engagement concentrated on the channels requiring informational navigation (specialist) and private expenditure (preventive dental care).

The second is the binary versus graduated structure of late-life fragility under the two regimes. The Italian typology splits the fragile stratum into two qualitatively distinct profiles, Fragile Resigned (13.2\%) and Fragile Depressed (26.7\%; jointly 39.9\%), spanning an active-depression-versus-settled-resignation contrast with separate policy responses, whereas the Swedish typology partitions the same population along a binary fragility line: a single Fragile cluster (11.5\%) against five autonomous profiles (88.5\%; \S\ref{sec:ch6-binary-graduated}). The universalistic regime absorbs the burden into a single targetable institutional segment; the familistic regime distributes it across multiple distinguishable household-level configurations.

The third is the universal digital saturation across the autonomous Swedish profiles against the resource-graduated adoption observed in Italy. Internet adoption ranges from 75\% to 99\% across the five autonomous Swedish profiles, with only the Fragile cluster below the 75\% line, while the Italian distribution ranges from 8.6\% to 74.8\% with no profile reaching it (\S\ref{sec:ch6-digital-cut}). The same diagnostic instrument cuts at the functional-fragility line in Sweden and at the resource-and-engagement gradient in Italy: the digital channel substitutes for the family-and-community infrastructure of healthcare and service navigation under the familistic regime, and complements the universal-access infrastructure under the universalistic regime.

The fourth is the age-driven rotation of the high-resource pole in Sweden against the parallel decline of both well-resourced Italian profiles. In Sweden the Connected Wealthy share contracts from 41\% (65--74) to 7\% (85+) while Asset Rich rises from 14\% to 31\%, crossing over in the 75--84 band (\S\ref{sec:ch6-age-rotation}); in Italy both Connected Active (31\% to 11\%) and Traditional Social (9\% to 3\%) decline in parallel with age. The rotation is the late-life trajectory permitted by a regime that absorbs the functional burden of advanced old age institutionally; the parallel decline is the trajectory where elderly care is delivered through intra-household and intra-family channels.

The four findings are not independent observations on a heterogeneous dataset. Each is the within-typology projection of the same regime contrast onto a different latent dimension of late-life experience: the structure of fragility, the operational meaning of the digital channel, the late-life trajectory of the high-resource pole, and the social-and-informational mediation of healthcare access. The thesis claim that the welfare regime is the structuring institution of the late-life partition is the joint reading of the four.

\section{Theoretical contribution}
\label{sec:ch9-theory}

The comparative welfare-regime literature initiated by \citet{espingandersen1990} and developed by \citet{ferrera1996}, \citet{anttonen1996}, \citet{bambra2007}, and \citet{naldini2008} establishes that the institutional treatment of social risk in the European welfare states sorts into a small number of regime types that produce systematically different aggregate outcomes on labour-market participation, family structure, healthcare access, and old-age income provision. The contribution of this thesis to that literature is to extend the regime reading from aggregate outcomes to the typological partition of the populations served by each regime: the Italian and Swedish over-65 populations do not only differ on aggregate indicators of late-life experience, they differ on the structural shape of the partition into archetypes that the same multidimensional segmentation pipeline returns from each population.

The methodological contribution to the segmentation literature is the inclusion of the subjective block as constituent input to the K-means partition. The empirical justification reported in \S\ref{sec:ch3-subjective-empirical} and formalised in the robustness battery of \S\ref{sec:ch7-cross-method} demonstrates that the inclusion of the subjective dimensions delivers a 41\% gain in external discriminative power on life satisfaction in Italy (and 22\% in Sweden) over the 4D specification that drops them, and preserves the qualitative contrast between Fragile Resigned and Fragile Depressed that the 4D specification collapses into a single fragile cluster. The position is that the segmentation literature on older adults should treat the subjective dimensions as constituent rather than derivative, in line with the gerontological tradition initiated by \citet{idlerbenyamini1997}, \citet{walker2005active}, and \citet{lalivedepinay2005}.

\section{Operational contribution to the silver-economy literature}
\label{sec:ch9-silver-economy}

The operational implications developed in \S\ref{sec:ch6-silver-economy} rest on three observations. First, addressability through direct-to-consumer digital products is structurally different across the two countries: in Sweden the autonomous profiles (88.5\% of the analytical sample) reach internet adoption of 75--99\% and are digitally addressable as a homogeneous block, against approximately 33\% of the Italian sample (Connected Active 25.4\% + Traditional Social 7.7\%) at comparable adoption. The Italian silver-economy market is structurally segmented into addressable and non-addressable cluster groups on dimensions that are themselves the signature of the regime.

Second, the unaddressed Italian Moderate Isolated segment (26.9\%) is the most direct case for product-design innovation: the binding constraint is not the absence of a digital channel but the absence of the social-and-informational mediation layer that would route the available channel towards healthcare and service utilisation. The candidate design is a digital channel combined with a social-mediation layer (peer-network platforms, community-based digital intermediaries, family-network-augmented service interfaces) that the Italian regime does not currently provide and that the Swedish typology does not require.

Third, the cross-country gap is structural rather than transitory: the Italian 65--74 cohort closing the digital gap to the Swedish reference under three rate scenarios takes approximately 6 to 30 years (\S\ref{sec:ch6-time-to-maturity}), and the cross-country dental-care gap is concentrated on the four matched pairs and, structurally, on the unmatched Moderate Isolated segment. The market-structure implication is therefore stated qualitatively: an addressable third of the Italian over-65 sample is served by direct-to-consumer and hybrid digital designs already targeted by the commercial silver-economy literature, against a non-addressable two thirds that would require the alternative design logics summarised in Table~\ref{tab:ch9-market-structure}.

\begin{table}[!htbp]
  \centering
  \small
  \caption[Italian typology by silver-economy addressability]{Italian over-65 typology by silver-economy addressability logic}
  \label{tab:ch9-market-structure}
  \renewcommand{\arraystretch}{1.2}
  \setlength{\tabcolsep}{4pt}
  \begin{tabularx}{\textwidth}{@{}>{\raggedright\arraybackslash\bfseries}p{2.4cm}>{\raggedright\arraybackslash}p{2.8cm}>{\centering\arraybackslash}p{1.2cm}>{\raggedright\arraybackslash}X@{}}
    \toprule
    \textbf{Addressability} & \textbf{Profile} & \textbf{Share \%} & \textbf{Operational design logic} \\
    \midrule
    \multicolumn{4}{@{}l}{\textit{Addressable through direct-to-consumer digital channel}} \\
    \midrule
    DTC digital      & Connected Active   & 25.4 & Direct-to-consumer digital products; high internet adoption (74.8\%). \\
    Hybrid intermediated & Traditional Social & 7.7  & Hybrid offline-digital designs; moderate internet (65.8\%) layered onto offline trust networks. \\
    \cmidrule(lr){2-4}
    \textbf{Subtotal addressable} & & \textbf{33.1} & Approximately one third of the analytical sample. \\
    \midrule
    \multicolumn{4}{@{}l}{\textit{Not addressable through direct-to-consumer digital channel}} \\
    \midrule
    Social-mediation & Moderate Isolated  & 26.9 & Social-and-informational mediation layer over digital channel; non-trivial internet (34.1\%) without dense informal networks. \\
    Affect-aware     & Fragile Depressed  & 26.7 & Affect-aware designs addressing the binding behavioural constraint (depressive affect, EURO-D 4.2); 17.3\% internet. \\
    Family-mediated  & Fragile Resigned   & 13.2 & Family or institutional intermediation; 8.6\% internet. Long-term-care welfare branch. \\
    \cmidrule(lr){2-4}
    \textbf{Subtotal non-addressable} & & \textbf{66.8} & Approximately two thirds of the analytical sample. \\
    \midrule
    \textbf{TOTAL} & & \textbf{99.9} & Italian over-65 analytical sample (\S\ref{sec:ch3-data}). \\
    \bottomrule
  \end{tabularx}
  \par\medskip
  \begin{minipage}{\linewidth}
    \footnotesize\textit{Note.} Profile shares from Table~\ref{tab:ch4-archetypes} and refer to the SHARE Wave~9 Italian analytical sample. Subtotals add to 99.9\% due to rounding. The table maps each typological segment to the operational design logic implied by its empirical signature on the digital and social-engagement dimensions (\S\ref{sec:ch4-profile-descriptions}, \S\ref{sec:ch6-addressability}); shares are descriptive of the sample partition (see \S\ref{sec:ch3-preprocessing}).
  \end{minipage}
\end{table}

The two-thirds non-addressable share is the substantive content of the Italian silver-economy design problem: a social-and-informational mediation layer for the Moderate Isolated segment (26.9\%), affect-aware designs for the Fragile Depressed segment (26.7\%), and family or institutional intermediation for the Fragile Resigned segment (13.2\%). The operational counterpart of this design problem, a working prototype of the missing mediation layer built as a B2B cohort-segmentation instrument over the typology, is developed in Chapter~\ref{ch:atlas}.

\section{Limitations}
\label{sec:ch9-limitations}

Three classes of limitations bound the substantive claims. The first is inferential and bears on the central claim: the regime reading rests on a cross-sectional, two-country comparison and cannot causally identify the welfare regime against the other respects in which Italy and Sweden differ (cohort composition, urbanisation, income level, cultural and historical context). The cross-wave evidence of \S\ref{sec:ch7-cross-wave} establishes structural persistence, not causal attribution; the regime interpretation is therefore the most parsimonious reading of the structural contrast rather than an identified effect, and generalisation beyond the two regimes requires the additional-country evidence set out in \S\ref{sec:ch9-future}. On the data side, the exclusion of proxy interviews biases the Fragile Resigned cluster downward in the prevalence of the most adverse configurations (severe cognitive impairment, institutional residence): the typology partitions the self-reporting over-65 cohort, not the full population. The single multiple-imputation replicate (\S\ref{sec:ch3-preprocessing}) introduces residual sensitivity at the Fragile--Moderate boundary, although item-level missingness is materially low (below 6\% on subjective indicators and below 9\% on income and wealth; \S\ref{sec:ch7-imputation}) and does not undermine the cluster centroids. On the modelling side, silhouette values lie below the 0.250 threshold (\S\ref{sec:ch7-internal-validity}), K-means/LCA convergence is moderate-to-substantial with the lowest rates on Wealthy Digital (49.3\%), Asset Rich (53.1\%) and Moderate Isolated (54.6\%) flagging geometrically thin cluster boundaries, and the 4D/5D ARI of 0.41 IT and 0.52 SE indicates only moderate partition agreement, so the subjective block materially reshapes the partition rather than perturbing it at the margins. The Hopkins statistic ($H = 0.75$ IT, $0.74$ SE) supports the existence of a real underlying structure, but the substantive findings of \S\ref{sec:ch9-findings} should be read as claims about cluster centroids and structural relations, not about marginal individual assignments.

\section{Future research}
\label{sec:ch9-future}

Five directions extend the analysis. First, applying the pipeline to additional SHARE countries (Germany, corporatist; Spain, conservative-Catholic; Poland and Czech Republic, post-socialist) would test whether the typological partition generalises into a regime-driven family of partitions \citep{bambra2007}. Second, building a transition matrix on the cross-wave panel intersection, with covariates entered as predictors of the transition probability, would identify the demographic and institutional drivers of late-life cluster trajectories and enable forecasting of cohort-level segment-share evolution. Third, a field-trial pairing existing digital service infrastructure with a structured peer-network or community-based intermediary, modelled on the documented evidence on social capital and healthcare navigation \citep{putnam2000}, would deliver an empirical test of the social-mediation-layer proposition for the Moderate Isolated segment. Fourth, integrating the cluster assignments with the within-Italy regional variation indices on healthcare access \citep{loscalzo2009} would yield a within-country counterpart of the cross-country comparison. Fifth, the Italian Silver Atlas platform of Chapter~\ref{ch:atlas} is currently supported analytically but untested in field; the two natural extensions are a qualitative comprehension study with Italian senior participants spanning the five clusters and a B2B field-trial with a healthcare practice and a pharmacy chain that observes how the cohort dashboard and action-card placeholders translate into operational workflows.

\section{Closing remarks}
\label{sec:ch9-closing}

The typological partition is the empirical instrument through which this thesis brings the welfare-regime contrast into the segmentation of late-life experience. The Italian and Swedish over-65 populations are partitioned into structurally different typologies by the same pipeline applied to the same input set, and the differences (four matched cross-country pairs and three regime-specific unmatched profiles, of which the Italian Moderate Isolated segment is the central country-specific finding) are the empirical signature of the welfare regime that surrounds each population.
